\documentclass[11pt,a4paper]{article}
\usepackage[utf8]{inputenc}
\usepackage[T1]{fontenc}
\usepackage[sc]{mathpazo}
\usepackage[scaled=0.92]{helvet}
\usepackage{microtype}
\usepackage[a4paper,top=2.6cm,bottom=2.4cm,left=2.4cm,right=2.4cm,headheight=14pt]{geometry}
\usepackage[table]{xcolor}
\usepackage{booktabs,array,enumitem,multirow}
\usepackage{tikz}
\usetikzlibrary{calc,positioning}
\usepackage{pgfplots}
\pgfplotsset{compat=1.16}
\usepackage[most]{tcolorbox}
\usepackage{titlesec}
\usepackage{needspace,etoolbox}
\usepackage{fancyhdr}
\usepackage{caption}
\usepackage{url}
\def\UrlBreaks{\do\/\do\-\do\_\do\.\do\=\do\?\do\&\do\a\do\b\do\c\do\d\do\e\do\f\do\g\do\h\do\i\do\j\do\k\do\l\do\m\do\n\do\o\do\p\do\q\do\r\do\s\do\t\do\u\do\v\do\w\do\x\do\y\do\z}
\Urlmuskip=0mu plus 1mu
\urlstyle{same}
\usepackage[hidelinks]{hyperref}

% ---- Palet (selaras dengan halaman dokumen pelengkap) ----
\definecolor{teal}{HTML}{0A6C6E}
\definecolor{tealdark}{HTML}{07484A}
\definecolor{ink}{HTML}{13262A}
\definecolor{muted}{HTML}{51666A}
\definecolor{rule}{HTML}{D3DEDC}
\definecolor{pale}{HTML}{EEF5F4}
\definecolor{amber}{HTML}{B7791F}
\definecolor{amberpale}{HTML}{FBF4E4}
\definecolor{well}{HTML}{CFE3E0}

\renewcommand{\tablename}{Tabel}
\renewcommand{\figurename}{Gambar}
\renewcommand{\refname}{Daftar Pustaka}
\color{ink}
\setlength{\parindent}{0pt}
\setlength{\parskip}{6.5pt plus 1pt}
\linespread{1.06}
\tolerance=1200
\emergencystretch=3em
\hyphenpenalty=1500
\setlist{itemsep=2.5pt,topsep=3pt,parsep=0pt,leftmargin=1.4em}
\setlist[itemize]{label={\color{teal}\textbullet}}
\renewcommand{\arraystretch}{1.22}
\setlength{\tabcolsep}{5pt}
\captionsetup{font={small,sf,color=muted},labelfont={bf,color=teal},labelsep=period,justification=raggedright,singlelinecheck=false,skip=5pt}

% ---- Judul bagian ----
\titleformat{\section}{\sffamily\bfseries\Large\color{ink}}{\color{teal}\thesection}{0.7em}{}[{\color{rule}\titlerule[0.6pt]}]
\titlespacing*{\section}{0pt}{18pt}{9pt}
\titleformat{\subsection}{\sffamily\bfseries\normalsize\color{tealdark}}{\color{teal}\thesubsection}{0.6em}{}
\titlespacing*{\subsection}{0pt}{11pt}{3pt}

% ---- Kepala dan kaki halaman ----
\pagestyle{fancy}
\fancyhf{}
\renewcommand{\headrulewidth}{0pt}
\fancyhead[L]{\sffamily\scriptsize\color{muted}\textsc{Alergen}\ \raisebox{0.15ex}{\tiny$\bullet$}\ Daftar Pustaka Terpadu}
\fancyhead[R]{\sffamily\scriptsize\color{muted}Kit Diagnostik Alergi Berbasis Patch}
\fancyfoot[L]{\sffamily\scriptsize\color{muted}Universitas Airlangga \ $\cdot$ \ IDS000011522}
\fancyfoot[R]{\sffamily\small\bfseries\color{teal}\thepage}
\renewcommand{\headrule}{\color{rule}\hrule width\headwidth height 0.5pt}
\renewcommand{\headrulewidth}{0.5pt}
\fancypagestyle{plain}{\fancyhf{}\renewcommand{\headrulewidth}{0pt}\fancyfoot[R]{\sffamily\small\bfseries\color{teal}\thepage}}

% ---- Tabel ----
\newcommand{\tH}[1]{\textcolor{white}{\sffamily\bfseries\small #1}}
\newcommand{\tabstart}{\small\rowcolors{2}{white}{pale}\arrayrulecolor{rule}}
\newcommand{\hdrrow}{\rowcolor{tealdark}}
\newcolumntype{L}[1]{>{\raggedright\arraybackslash}p{#1}}
\newcolumntype{R}[1]{>{\raggedleft\arraybackslash}p{#1}}
\setlength{\aboverulesep}{0pt}\setlength{\belowrulesep}{0pt}
\renewcommand{\toprule}{\arrayrulecolor{tealdark}\specialrule{0.9pt}{0pt}{0pt}\arrayrulecolor{rule}}
\renewcommand{\bottomrule}{\arrayrulecolor{tealdark}\specialrule{0.9pt}{0pt}{0pt}\arrayrulecolor{rule}}

% ---- Kotak catatan dan kartu angka ----
\newtcolorbox{catatan}[1][Catatan]{enhanced,breakable,colback=amberpale,colframe=amberpale,boxrule=0pt,
  borderline west={2.2pt}{0pt}{amber},arc=1.5pt,left=10pt,right=8pt,top=6pt,bottom=6pt,
  before skip=8pt,after skip=10pt,
  title={\sffamily\bfseries\small\color{amber}\MakeUppercase{#1}},coltitle=amber,
  attach boxed title to top left={yshift=-2mm,xshift=0pt},boxed title style={empty,left=0pt,top=0pt,bottom=0pt},
  fonttitle=\sffamily,coltext=ink,before upper={\vspace{2pt}}}
\newtcolorbox{kartu}[2]{enhanced,colback=pale,colframe=pale,boxrule=0pt,arc=3pt,
  left=8pt,right=6pt,top=6pt,bottom=6pt,halign=left,
  borderline west={2.5pt}{0pt}{teal},
  before upper={{\sffamily\bfseries\LARGE\color{teal}#1}\par\vspace{1pt}},
  fontupper=\sffamily\scriptsize\color{muted},height=1.95cm,valign=top}
\newcommand{\kartuRow}[4]{%
 \begin{tcbraster}[raster columns=#1,raster equal height,raster column skip=7pt,raster row skip=7pt]#2#3#4\end{tcbraster}\vspace{9pt}}

% ---- Sampul ----
\newcommand{\patchmotif}[2]{% #1 skala, #2 warna garis
 \begin{scope}[scale=#1]
 \draw[#2,line width=1.2pt,rounded corners=9pt] (0,0) rectangle (15,6);
 \foreach \r in {0,1}{\foreach \c in {0,...,4}{
   \draw[#2,line width=1pt] ({1.9+\c*2.7},{4.1-\r*2.2}) circle (0.9);}}
 \end{scope}}

\newcommand{\makecover}[5]{% kode, judul, subjudul, tagline, ringkasan satu kalimat
\begin{titlepage}
\thispagestyle{empty}
\begin{tikzpicture}[remember picture,overlay]
 \fill[tealdark] (current page.north west) rectangle ($(current page.north east)+(0,-17.2cm)$);
 \fill[teal!85!tealdark] ($(current page.north east)+(0,-17.2cm)$) -- ($(current page.north east)+(-10cm,-17.2cm)$) -- ($(current page.north east)+(0,-10.8cm)$) -- cycle;
 \begin{scope}[shift={($(current page.north east)+(-9.9cm,-6.4cm)$)}]
   \patchmotif{0.5}{white!35}
 \end{scope}
 \node[anchor=north west,text=white!70,font=\sffamily\small\bfseries,align=left] at ($(current page.north west)+(2.4cm,-2.7cm)$) {\textsc{Alergen \ $\cdot$ \ Universitas Airlangga}};
 \node[anchor=north west,text=white,font=\sffamily\bfseries\fontsize{44}{44}\selectfont] at ($(current page.north west)+(2.4cm,-7.4cm)$) {#1};
 \node[anchor=north west,text=white,font=\sffamily\bfseries\fontsize{26}{31}\selectfont,text width=14.2cm,align=left] at ($(current page.north west)+(2.4cm,-10.2cm)$) {#2};
 \node[anchor=north west,text=white!85,font=\sffamily\fontsize{13}{17}\selectfont,text width=14cm,align=left] at ($(current page.north west)+(2.4cm,-14.1cm)$) {\hyphenpenalty=10000\exhyphenpenalty=10000 #3};
 \node[anchor=north west,text=ink,font=\fontsize{12}{17}\selectfont,text width=15.2cm,align=left] at ($(current page.north west)+(2.4cm,-18.8cm)$) {#4};
 \draw[teal,line width=2pt] ($(current page.north west)+(2.4cm,-21.9cm)$) -- ++(2.2cm,0);
 \node[anchor=north west,text=muted,font=\sffamily\small,text width=15cm,align=left] at ($(current page.north west)+(2.4cm,-22.4cm)$) {%
  \textbf{\color{ink}Kit Diagnostik Alergi Berbasis Patch (Alergen)}\\[3pt]
  Universitas Airlangga, Surabaya\\
  Paten Sederhana IDS000011522\\
  Oktober 2026};
 \node[anchor=south east,text=muted,font=\sffamily\scriptsize,align=right] at ($(current page.south east)+(-2.4cm,1.6cm)$) {#5};
\end{tikzpicture}
\end{titlepage}
\newpage}

\definecolor{sr1}{HTML}{8AA6AB}
\definecolor{sr2}{HTML}{3FA7A3}
\definecolor{sr3}{HTML}{0A6C6E}
\definecolor{sr4}{HTML}{B7791F}
\pgfplotsset{natbase/.style={width=0.95\linewidth,height=6.2cm,
  every axis label/.append style={font=\sffamily\scriptsize,color=muted},
  tick label style={font=\sffamily\scriptsize,color=muted},
  legend style={font=\sffamily\scriptsize,draw=rule,fill=white,at={(0.5,-0.2)},anchor=north},
  grid=major,grid style={color=rule,line width=0.4pt},
  axis line style={color=rule},tick style={color=rule},
  scaled ticks=false,/pgf/number format/.cd,1000 sep={.},fixed,precision=0,/tikz/.cd},
  natplot/.style={natbase,xtick={1,2,3,4,5},xticklabels={2028,2029,2030,2031,2032}}}
\hypersetup{pdftitle={Daftar Pustaka Terpadu},pdfauthor={Tim Alergen, Universitas Airlangga}}

\pretocmd{\section}{\Needspace{7\baselineskip}}{}{}
\pretocmd{\subsection}{\Needspace{4\baselineskip}}{}{}

\usepackage{float}
\begin{document}
\makecover{Pustaka}{Daftar Pustaka Terpadu}{Rujukan yang mendasari studi kelayakan, dampak ekonomi dan sosial, CoA, dan dokumen produksi patch Alergen}{Dokumen ini menghimpun seluruh rujukan dari empat dokumen pendukung dalam satu daftar, dikelompokkan menurut topik, dengan petunjuk dokumen yang memakainya. Tiap rujukan hanya muncul satu kali.}{39 rujukan, tujuh kelompok topik}

\section*{Ringkasan}
\addcontentsline{toc}{section}{Ringkasan}

\kartuRow{4}
{\begin{kartu}{39}{}rujukan unik pada daftar terpadu\end{kartu}}
{\begin{kartu}{7}{}kelompok topik, dari epidemiologi sampai metode ekstrak\end{kartu}}
{\begin{kartu}{4}{}dokumen pendukung yang dirujuk\end{kartu}
\begin{kartu}{2026}{}tanggal akses terakhir semua laman, Oktober\end{kartu}}

Daftar ini memuat rujukan yang dikutip pada empat dokumen berikut. Rujukan yang dipakai lebih dari satu dokumen ditulis sekali, dan kolom ``Dipakai pada'' menunjukkan dokumen yang mengutipnya.

\begin{table}[H]
\centering
\caption{Dokumen yang dirujuk.}
\tabstart
\begin{tabular}{@{}L{5.8cm}L{9.4cm}@{}}
\toprule
\hdrrow\tH{Dokumen} & \tH{Isi} \\
\midrule
Studi Kelayakan Skala Nasional & Pasar, strategi peluncuran, produksi, keuangan, kekayaan intelektual, risiko \\
Dampak Ekonomi dan Sosial Skala Nasional & Beban penyakit, dampak bagi pasien, fasilitas, industri, dan sistem kesehatan \\
Certificate of Analysis (CoA) & Sertifikat analisis 10 ekstrak alergen dan patch \\
Dokumen Produksi & Proses, bahan, catatan batch, dan kestabilan \\
\bottomrule
\end{tabular}
\end{table}

Laman web dicantumkan dengan tanggal akses atau tanggal terbit bila tersedia. Angka dari sumber sekunder (siaran pers, ringkasan, atau pemberitaan) diberi keterangan sumbernya pada rujukan yang bersangkutan.

\Needspace{44\baselineskip}
\section{Beban penyakit alergi dan epidemiologi}

\begin{table}[H]
\centering
\caption{Rujukan: beban penyakit alergi dan epidemiologi.}
{\footnotesize\tabstart\setlength{\tabcolsep}{3.5pt}
\begin{tabular}{@{}rL{10.3cm}L{2.8cm}@{}}
\toprule
\hdrrow\tH{No} & \tH{Rujukan} & \tH{Dipakai pada} \\
\midrule
1 & World Health Organization. Asthma, fact sheet, 28 April 2026. \url{https://www.who.int/news-room/fact-sheets/detail/asthma} & Dampak \\
2 & Kementerian Kesehatan RI. Riskesdas 2018 (asma terdiagnosis dokter, 2,4\%) dan Infodatin Penderita Asma di Indonesia, 2019. Dikutip melalui sumber sekunder. & Dampak \\
3 & ANTARA. Prevalensi alergi susu sapi pada anak Indonesia 2 sampai 7,5\% (data Ikatan Dokter Anak Indonesia). 25 Juni 2024. & Dampak \\
4 & Soegiarto G, Abdullah MS, Damayanti LA, Suseno A, Effendi C. The prevalence of allergic diseases in school children of metropolitan city in Indonesia shows a similar pattern to that of developed countries. Asia Pac Allergy. 2019;9(2):e17. doi:10.5415/apallergy.2019.9.e17 & Dampak \\
5 & Norb\"ack D, Hashim JH, Hashim Z, et al. Fractional exhaled nitric oxide (FeNO) among school children in Java and Sumatra, Indonesia: associations with respiratory symptoms, house dust mite sensitization and the home environment. J Asthma. 2024;61(12):1772-1780. doi:10.1080/02770903.2024.2383627 & Dampak \\
6 & Zuberbier T, L\"otvall J, Simoens S, Subramanian SV, Church MK. Economic burden of inadequate management of allergic diseases in the European Union: a GA2LEN review. Allergy. 2014;69(10):1275-1279. doi:10.1111/all.12470. Dikutip dari siaran pers Wiley: \url{https://eurekalert.org/news-releases/494816} & Dampak \\
7 & Gupta R, Holdford D, Bilaver L, et al. The economic impact of childhood food allergy in the United States. JAMA Pediatr. 2013;167(11):1026-1031. doi:10.1001/jamapediatrics.2013.2376. Dikutip dari ringkasan: \url{https://www.newswise.com/articles/study-estimates-economic-impact-of-childhood-food-allergies} & Dampak \\
8 & Cay M, Onal C, Herbert L, et al. Psychosocial burden of multiple IgE-mediated food allergies in pediatric patients and caregivers in the FORWARD study. Nutrients. 2026;18(11):1745. doi:10.3390/nu18111745 & Dampak \\
9 & Wai CYY, Santos AF, Leung TF, Sampson HA. Novel strategies for food allergy diagnosis. J Allergy Clin Immunol Pract. 2026;14(4):751-760. doi:10.1016/j.jaip.2025.11.022 & Dampak \\
\bottomrule
\end{tabular}}
\end{table}

\Needspace{28\baselineskip}
\section{Demografi, layanan kesehatan, dan jaminan kesehatan}

\begin{table}[H]
\centering
\caption{Rujukan: demografi, layanan kesehatan, dan jaminan kesehatan.}
{\footnotesize\tabstart\setlength{\tabcolsep}{3.5pt}
\begin{tabular}{@{}rL{10.3cm}L{2.8cm}@{}}
\toprule
\hdrrow\tH{No} & \tH{Rujukan} & \tH{Dipakai pada} \\
\midrule
10 & BPS. Hasil Survei Penduduk Antar Sensus (SUPAS) 2025: 284,67 juta jiwa. Dikutip melalui Bisnis.com, 5 Mei 2026. \url{https://kabar24.bisnis.com/read/20260505/15/1971460/sensus-bps-jumlah-penduduk-ri-capai-28467-juta-jiwa-2025-tumbuh-melambat} & Kelayakan \\
11 & Sekretariat Kabinet RI. Hasil Sensus Penduduk 2020; BPS: Meski lambat, ada pergeseran penduduk antarpulau. \url{https://setkab.go.id/hasil-sensus-penduduk-2020-bps-meski-lambat-ada-pergeseran-penduduk-antarpulau/} & Kelayakan, Dampak \\
12 & Bisnis.com. BPJS Kesehatan klaim kerja sama dengan 3.120 RS pada 2023. 14 Januari 2024. \url{https://finansial.bisnis.com/read/20240114/215/1732195/bpjs-kesehatan-klaim-kerja-sama-dengan-3120-rs-pada-2023} & Kelayakan, Dampak \\
13 & Warta Ekonomi. Cakupan JKN tembus 98 persen penduduk Indonesia, BPJS catat peserta capai 282,7 juta jiwa. 27 Januari 2026. \url{https://wartaekonomi.co.id/read598290/cakupan-jkn-tembus-98-penduduk-indonesia-bpjs-catat-peserta-capai-2827-juta-jiwa} & Kelayakan, Dampak \\
14 & Kabarbursa.com. Menkes: RI masih butuh 29.000 dokter spesialis. 6 Mei 2024. \url{https://www.kabarbursa.com/market-hari-ini/menkes-ri-masih-butuh-29-000-dokter-spesialis} & Dampak \\
\bottomrule
\end{tabular}}
\end{table}

\Needspace{44\baselineskip}
\section{Harga, pasar, dan keuangan}

\begin{table}[H]
\centering
\caption{Rujukan: harga, pasar, dan keuangan.}
{\footnotesize\tabstart\setlength{\tabcolsep}{3.5pt}
\begin{tabular}{@{}rL{10.3cm}L{2.8cm}@{}}
\toprule
\hdrrow\tH{No} & \tH{Rujukan} & \tH{Dipakai pada} \\
\midrule
15 & Global Market Insights. Allergy Diagnostics Market. Diakses Oktober 2026. \url{https://www.gminsights.com/industry-analysis/allergy-diagnostics-market} & Kelayakan \\
16 & ResearchAndMarkets.com. Global Allergy Diagnostics Market Report 2025. GlobeNewswire, 1 Mei 2025. \url{https://www.globenewswire.com/news-release/2025/05/01/3072134/28124/en/Global-Allergy-Diagnostics-Market-Report-2025-Forecast-to-Reach-13-4-Billion-by-2034-Reimbursement-Support-and-Technological-Innovation-Drive-Growth-at-a-CAGR-of-9-20.html} & Kelayakan \\
17 & Alodokter. Tes Alergi, daftar tarif fasilitas di Surabaya. Diakses Oktober 2026. \url{https://www.alodokter.com/cari-rumah-sakit/kulit-dan-kelamin/tes-alergi/surabaya} & Kelayakan, Dampak \\
18 & Rumah Sakit Royal Progress. Paket Cek Alergi, IgE Atopy 54 Alergen (Rp2.200.000, termasuk konsultasi dokter). Diakses Oktober 2026. \url{https://oncare.royalprogress.com/en/s/rumah-sakit-royal-progress/detail-produk/paket-cek-alergi-ige-54-alergen} & Kelayakan, Dampak \\
19 & Medicastore. Tes Alergi (uji tusuk kulit 33 alergen Rp350.000 sampai Rp600.000; IgE spesifik Rp200.000 sampai Rp300.000 per alergen). 9 Januari 2009. \url{https://medicastore.com/artikel/1964/tes-alergi} & Kelayakan, Dampak \\
20 & Laboratorium Klinik Pramita. Allergy Test Panels (IgE spesifik 96 alergen; harga tidak dicantumkan). 2 Oktober 2025. \url{https://pramita.co.id/en/inspiration/allergy-test-panels} & Kelayakan, Dampak \\
21 & Laboratorium Klinik Pramita. Food Allergy Check Panel (IgE spesifik 54 jenis; IgG Food Tolerance 46 jenis makanan; harga tidak dicantumkan). Diakses Oktober 2026. \url{https://pramita.co.id/en/inspiration/food-allergy-check-panel} & Kelayakan, Dampak \\
22 & Daftar harga tes alergi pada Siloam Hospitals, Tirta Medical Centre, RS Mitra Keluarga, Klinik Utama DR Indrajana, dan Laboratorium Klinik Pramita. Dihimpun Oktober 2026. & Kelayakan, Dampak \\
23 & Cermati Invest. Weekly Update 19 Oktober 2026. \url{https://invest.cermati.com/artikel/amp/cermati-invest-weekly-update-19-oktober-2026} & Kelayakan \\
\bottomrule
\end{tabular}}
\end{table}

\Needspace{32\baselineskip}
\section{Regulasi, pedoman mutu, dan produk pembanding}

\begin{table}[H]
\centering
\caption{Rujukan: regulasi, pedoman mutu, dan produk pembanding.}
{\footnotesize\tabstart\setlength{\tabcolsep}{3.5pt}
\begin{tabular}{@{}rL{10.3cm}L{2.8cm}@{}}
\toprule
\hdrrow\tH{No} & \tH{Rujukan} & \tH{Dipakai pada} \\
\midrule
24 & European Medicines Agency. Guideline on allergen products: production and quality issues. EMEA/CHMP/BWP/304831/2007, disahkan 20 November 2008, berlaku 20 Mei 2009. \url{https://www.ema.europa.eu/en/allergen-products-production-quality-issues-scientific-guideline} & Kelayakan, CoA, Produksi \\
25 & European Pharmacopoeia. Bab 5.1.4, Microbiological quality of non-sterile pharmaceutical preparations and substances for pharmaceutical use; bab 2.6.12, 2.6.13, dan 2.6.14. & CoA, Produksi \\
26 & International Council for Harmonisation. ICH Q1A(R2): Stability testing of new drug substances and products, 2003; ICH Q5C: Stability testing of biotechnological/biological products, 1995. & CoA, Produksi \\
27 & Hukumonline Pro. Drug registration procedures updated: BPOM releases list of 11 biological products and expands evaluation pathways. 19 Agustus 2025 (ringkasan terbuka). \url{https://pro.hukumonline.com/a/lt68a3dbcf7be8e/drug-registration-procedures-updated--bpom-releases-list-of-11-biological-products-and-expands-evaluation-pathways/} & Kelayakan \\
28 & Veritask. Peraturan Pemerintah Nomor 15 Tahun 2026 mengubah struktur tarif registrasi, inspeksi, dan sertifikasi pada BPOM. Diakses Oktober 2026. \url{https://veritask.ai/id/artikel/peraturan-pemerintah-nomor-15-tahun-2026-mengubah-struktur-tarif-registrasi-inspeksi-dan-sertifikasi-pada-badan-pengawas-obat-dan-makanan} & Kelayakan \\
29 & FDA, Center for Biologics Evaluation and Research. Safety and utilization review for the Pediatric Advisory Committee (T.R.U.E. TEST, STN 103738), September 2021. \url{https://www.fda.gov/media/151565/download} & Kelayakan \\
\bottomrule
\end{tabular}}
\end{table}

\Needspace{24\baselineskip}
\section{Industri dalam negeri dan kebijakan produk lokal}

\begin{table}[H]
\centering
\caption{Rujukan: industri dalam negeri dan kebijakan produk lokal.}
{\footnotesize\tabstart\setlength{\tabcolsep}{3.5pt}
\begin{tabular}{@{}rL{10.3cm}L{2.8cm}@{}}
\toprule
\hdrrow\tH{No} & \tH{Rujukan} & \tH{Dipakai pada} \\
\midrule
30 & Bio Farma. Bio Farma segera luncurkan pendeteksi diabetes mellitus tipe I. Pengumuman Bio Farma. \url{https://www.biofarma.co.id/en/announcement/detail/bio-farma-segera-luncurkan-pendeteksi-diabetes-mellitus-type-i} & Kelayakan \\
31 & Kabarbursa.com. Pemerintah genjot produk lokal masuk e-katalog LKPP. 16 Oktober 2025. \url{https://www.kabarbursa.com/makro/pemerintah-genjot-produk-lokal-masuk-e-katalog-lkpp} & Kelayakan, Dampak \\
32 & BCA Sekuritas. Kemenperin perkuat pasar industri alat kesehatan lewat P3DN. 8 September 2025. \url{https://bcasekuritas.co.id/latest-news/news/kemenperin-perkuat-pasar-industri-alat-kesehatan-lewat-p3dn} & Kelayakan, Dampak \\
33 & Kementerian Kesehatan RI. Pemerintah targetkan 60 persen produksi alkes gunakan komponen lokal. 27 Agustus 2022. \url{https://www.kemkes.go.id/id/pemerintah-targetkan-60-persen-produksi-alkes-gunakan-komponen-lokal} & Kelayakan, Dampak \\
\bottomrule
\end{tabular}}
\end{table}

\Needspace{24\baselineskip}
\section{Kekayaan intelektual dan lisensi}

\begin{table}[H]
\centering
\caption{Rujukan: kekayaan intelektual dan lisensi.}
{\footnotesize\tabstart\setlength{\tabcolsep}{3.5pt}
\begin{tabular}{@{}rL{10.3cm}L{2.8cm}@{}}
\toprule
\hdrrow\tH{No} & \tH{Rujukan} & \tH{Dipakai pada} \\
\midrule
34 & Izin.co.id. Jangka waktu paten di Indonesia. Diakses Oktober 2026. \url{https://izin.co.id/blog/jangka-waktu-paten/} & Kelayakan \\
35 & Gunderson Dettmer. University License Presentation, 7 Februari 2024. \url{https://www.gunder.com/content/uploads/2024/02/Gunderson-Dettmer-University-License-Presentation-2-7-24.pdf} & Kelayakan \\
36 & BVWire. Royalty rates for university-bred technology. Isu 186-3, 21 Maret 2018. \url{https://www.bvresources.com/articles/bvwire/royalty-rates-for-university-bred-technology} & Kelayakan \\
37 & Licensing Executives Society International. Royalty rates and license fees for technology. \url{https://lesi.org/wp-content/uploads/2024/04/royalty-rates-license-fees-for-technology.pdf} & Kelayakan \\
\bottomrule
\end{tabular}}
\end{table}

\Needspace{16\baselineskip}
\section{Bahan dan metode ekstrak alergen}

\begin{table}[H]
\centering
\caption{Rujukan: bahan dan metode ekstrak alergen.}
{\footnotesize\tabstart\setlength{\tabcolsep}{3.5pt}
\begin{tabular}{@{}rL{10.3cm}L{2.8cm}@{}}
\toprule
\hdrrow\tH{No} & \tH{Rujukan} & \tH{Dipakai pada} \\
\midrule
38 & Pratama YA, Marhaeny HD, Kasatu SM, Rohmah L, Nurhan AD, Rahmadi M, Khotib J. The optimisation and standardisation of Indonesian house dust mites allergenic extract as a desensitising agent. Pharmacy Education 2023;23(4):137-139. doi:10.46542/pe.2023.234.137139 & CoA, Produksi \\
39 & Marhaeny HD, Rohmah L, Pratama YA, Kasatu SM, Miatmoko A, Addimaysqi R, dkk. Shrimp allergen extract immunotherapy induces prolonged immune tolerance in a gastro-food allergy mouse model. PLoS ONE 2024;19(12):e0315312. doi:10.1371/journal.pone.0315312 & CoA, Produksi \\
\bottomrule
\end{tabular}}
\end{table}


\end{document}
